\begin{myChapterIntro}{本章涵盖}
\begin{itemize}
\item 应用的功能性需求与非功能性需求
\item 什么是好的需求，以及如何获取需求
\item 用例
\item 功能规格说明
\item 分析需求以得到应用的初始类
\end{itemize}
\end{myChapterIntro}

在我们开始操心如何\emph{正确地构建应用}（把应用设计好）之前，必须先确保我们要构建的是\emph{正确的应用}。一个不满足客户需求的应用，无论设计得多好，都是失败的、不成功的应用。应用的\emph{客户}可以是未来的最终用户（包括你自己）、要求你编写该应用的经理、聘用你担任软件顾问或承包商的人，或者任何希望应用成功的利益相关者。

\myGraphic{0.980}{images/CH03_00_UN01_Mak.png}{}

本章将学习如何为应用获取好的需求。然后，我们将了解如何分析这些需求，以得到初始的类集合。请记住，需求会发生变化，我们必须相应地设计，并采用迭代式开发。

\section{应用设计的序曲}

图 3.1 展示了应用开发过程中主要活动与里程碑的时间线。用粗体轮廓标出的活动（本章会介绍）是应用设计至关重要的序曲。

\myGraphic{0.891}{images/CH03_F01_Mak.png}{图 3.1 开发应用时主要活动的时间线。用粗体轮廓标出的活动是应用设计至关重要的前奏。在获取需求（Get Requirements）和功能规格说明（Functional Specification）之后用虚线表示的时间段表明，它们的冻结日期必须是灵活的。}

假设客户请我们开发一个比第 2 章的图书目录应用更高级的版本。客户对“图书目录服务器”有一些想法，希望它把图书数据存储在某种数据库中，并让查找图书的图书馆顾客通过一个表单界面填写并提交想要查询的图书属性。此外，还有一种扮演图书管理员角色的特殊最终用户，负责向目录中添加新书，以及更新或删除已有的书（图 3.2）。我们如何把这些模糊的想法转化为确定的需求，以确保最终开发出的应用能够完全满足客户？需求分为两类：功能性需求和非功能性需求。

\myGraphic{0.587}{images/CH03_F02_Mak.png}{图 3.2 更高级版本的图书目录应用。图书管理员与图书目录服务器交互，向目录中输入新书并更新和删除已有的书。顾客与服务器交互以查找图书。图书存储在后端数据库中。}

\section{功能性需求：应用必须做什么？}

\emph{功能性需求}（functional requirement）规定了一个成功应用必须执行的某项具体操作，或者必须允许最终用户执行的某项具体操作。在第 2 章中，我们看到了示例图书目录应用的功能性需求。基于这些需求，应用必须执行的操作清单包括

\begin{itemize}
\item 存储所有图书的属性，包括书名以及作者的名和姓
\item 为小说类图书存储年份和体裁属性，为烹饪书存储地区属性，为实用指南类图书存储主题属性
\item 执行不区分大小写的字符串匹配
\item 允许在图书搜索中指定不关心的目标属性
\end{itemize}

基于这些功能性需求，应用必须允许最终用户执行的操作清单包括

\begin{itemize}
\item 向目录中添加小说、烹饪书和实用指南类图书
\item 在目录中查找与给定的一组目标属性匹配的图书
\end{itemize}

我们可以用 \emph{must}、\emph{shall} 等语气强烈的情态动词把这些操作改写为需求：

\begin{itemize}
\item 用户\emph{必须}能够向目录中添加小说、烹饪书和实用指南类图书及其属性。
\item 每本图书的属性\emph{应}为书名以及作者的姓和名。
\item 小说的属性\emph{应}包括出版年份和体裁：冒险、经典、侦探、奇幻、历史、恐怖、言情或科幻。
\item 烹饪书的属性\emph{应}包括其地区：中国、法国、印度、意大利、墨西哥、伊朗或美国。
\item 实用指南类图书的属性\emph{应}包括其主题：素描、绘画或写作。
\item 用户\emph{必须}能够按给定的一组目标属性在目录中查找匹配的图书。
\item 图书搜索时匹配属性的字符串比较\emph{应}不区分大小写。
\item 用户\emph{必须}能够为图书搜索指定不关心的属性。
\end{itemize}

当我们倾听客户的诉求并将其愿望转化为书面需求时，不应使用 \emph{should} 或 \emph{could} 这类语气较弱的情态动词来表述。

\myGraphic{0.764}{images/CH03_F02_UN02_Mak.png}{}

\section{非功能性需求：对应用的约束}

\emph{非功能性需求}（nonfunctional requirement）规定了应用为取得成功必须满足的某项具体限制或约束。这类需求涉及性能、平台、可维护性等问题。

在第 2 章中，我们没有为图书目录应用列出非功能性需求。下面是该应用的一组非功能性需求：

\begin{itemize}
\item 一次图书搜索必须在 2 秒以内完成。
\item 应用必须能在 Windows、macOS 和 Linux 平台上运行。
\item 用户界面应与此前版本的 UI 相似。
\item 显示的消息应可定制为英语、西班牙语或越南语。
\end{itemize}

\myGraphic{0.811}{images/CH03_F02_UN03_Mak.png}{}

不要以为非功能性需求不如功能性需求重要！如果性能无法接受，或者无法在用户的平台上运行，那么一个设计良好、功能完备的应用同样是失败的。

\begin{mySidebar}{国际化说明}

上面列出的最后一条非功能性需求涉及\emph{国际化}（internationalization），有时缩写为 \emph{I18N}（因为单词开头的 \emph{I} 和结尾的 \emph{N} 之间正好有 18 个字母）。把应用适配到特定区域（locale）的过程称为\emph{本地化}（\emph{L10N}），它涉及更改文本消息所用的自然语言（英语、西班牙语等），以及进行日期、时间、货币格式、字符编码标准等其他文本更改。如果要开发国际化的应用，当区域（locale）变化时，应把应用中受影响的部分封装起来。
\end{mySidebar}

需求不应规定如何设计应用来满足该需求。例如，关于用户界面的那条非功能性需求并没有说要把旧 UI 的表单纳入新 UI。需求不应包含设计和实现细节。

\myGraphic{0.893}{images/CH03_F02_UN04_Mak.png}{}

\section{什么是好的需求？}

怎样才算是一组理想的需求？用于开发设计良好的应用的理想需求，本身就应满足以下标准：

\begin{itemize}
\item \emph{清晰性}——需求必须用非技术性、不含术语的语言清晰地编写，以便客户和软件开发者都能理解。
\item \emph{一致性}——需求之间不能相互矛盾。例如，不能有一条需求规定应用必须能在 Windows 和 macOS 上运行，另一条需求又规定应用只能运行在 Linux 上。
\item \emph{正确性}——每条需求都必须正确。例如，某个医疗应用的一条需求规定该应用必须处理男性患者的妊娠数据，这显然是错误的。
\item \emph{完整性}——需求中的任何缺漏都可能导致软件开发者做出错误猜测，进而使应用无法满足客户的需求。
\item \emph{切合实际}——不要包含无法满足的需求，例如过于乐观的性能指标。
\item \emph{可验证性}——能否通过测试应用来确保它满足每一条需求？
\item \emph{可追溯性}——我们能否把每条需求追溯到应用的某项功能或约束？我们不希望应用漏掉任何需求。反过来，我们能否把每项功能或约束追溯到某条需求？我们不希望给应用加上多余的功能。
\end{itemize}

\myGraphic{0.686}{images/CH03_F02_UN05_Mak.png}{}

\section{如何获取需求}

我们如何获取应用的需求？需求从何而来？既然需求决定了应用怎样才算成功，那么一个显而易见的起点就是向客户询问这个应用。这可能涉及一系列长时间的访谈——除非应用微不足道，否则这几乎肯定不是一次性的访谈。要提问并求得澄清：这是你想要的效果吗？你那样说是什么意思？你希望应用做这个还是做那个？

这种收集需求的访谈需要良好的人际交往能力。我们要去发现客户想要什么，而不是把自己的偏见带入其中。我们要跨越一道巨大的鸿沟——我们假定客户是其所在领域的专家（比如银行应用的金融领域），而我们才是软件开发方面的专家。客户的专业知识和词汇可能与我们的很少有共同之处。

但有一个大陷阱：\emph{客户并不总是知道他们想要什么！}不要施加太大压力逼他们拿出需求，否则他们可能会开始胡乱编造。你的客户可能在想：“我付给这些开发者这么多钱，他们还有截止日期。我不给他们更多需求，他们就动不起来。我还是给他们\emph{点东西}，好让他们开始写代码吧。”这必定会导致糟糕的需求，而糟糕的需求只会带来糟糕的应用。

因此，软件开发者必须努力通过其他途径获取需求。一种方法是观察人们当前完成任务的方式，并设想应用如何改进他们的工作表现。观察那些靠手工或使用遗留软件完成的具体任务，并构思出让应用更好地自动完成这些任务的需求。

在一些软件组织中，项目经理等比程序员更擅长沟通的团队成员负责访谈客户，以取得初步的需求清单。不过，开发者仍需与他们要解决的问题有一定程度的实际接触，这依然很重要。

开发团队成员之间可能需要召开多次会议，才能最终确定功能性需求和非功能性需求的清单。在理想情况下，应用的客户会成为开发团队的一名虚拟成员，随时可以解答关于需求的问题。若无法做到这一点，那么保持开发者与客户之间的沟通渠道尽可能畅通就至关重要。客户与开发团队之间的反馈回路必须紧密，并始终处于活跃状态。

\subsection{3.5.1 一个简短的需求案例研究}

图 3.3 展示了一个简单的场景。公共图书馆过去使用卡片目录，那是一排排实体的抽屉，里面存放着记录图书馆藏书属性的索引卡。卡片按书名、作者姓氏和主题分别按字母顺序排列。因此，每本书至少需要三张卡片。每个抽屉都按其所存放卡片的字母范围贴上标签。想要找书的图书馆顾客会查阅卡片目录，看图书馆是否有这本书；如果有，卡片会指明在哪里能找到它。图书管理员负责维护卡片目录。无论对图书管理员还是顾客而言，这都是一个非常费时的手工过程。因此，我们的客户——图书馆馆长——聘我们来将目录自动化。

\myGraphic{0.886}{images/CH03_F03_Mak.png}{图 3.3 要用软件应用将图书馆老式的卡片目录自动化，必须先获取需求，才能开发出正确的应用。}

我们必须先从客户那里获取需求，才能开发出正确的应用。我们访谈图书馆馆长，了解管理员们如何维护卡片目录，并询问他们认为我们能做些什么来改进他们的工作。我们还访谈图书馆的顾客，了解他们如何在卡片目录中查找图书，并询问我们能做些什么来改进他们的搜索。但也许管理员只能对自身工作提出渐进式的改进建议，比如用一个在线表单来制作并打印抽屉用的卡片。同样，图书馆顾客可能也缺乏想象力，提不出更好的检索方法。因此，我们必须观察管理员的工作和顾客的检索过程，并运用我们在计算机应用方面的经验，构思出更多的需求。我们可以制作并展示新应用的原型，以激发更好的需求。获取需求也是一个迭代式的过程。

\myGraphic{0.780}{images/CH03_F03_UN06_Mak.png}{}

\subsection{3.5.2 明示需求与隐含需求}

管理员和顾客提供的是\emph{明示}的应用需求。通过领会言外之意，我们可以构想出\emph{隐含}的需求。没错，我们的受访者说要需求 A 和 C。但凭我们自己的经验（或常识）可知，他们还需要需求 B，才能真正把他们的工作自动化。

在图 3.1 的时间线上，客户和开发者必须就“需求冻结”（Requirements freeze）这一日期达成一致。按理说，该日期之后需求就不会再变动了。但收集到的需求往往并不理想，因此我们应当预料到它们可能仍会变化，也可能还会有额外的需求。不过，需求必须在冻结日期之后不久就稳定下来，以便软件设计能够趋于稳定，并最终按时完成应用。好的应用设计会封装应用中那些最有可能受需求变化影响的部分。设计良好的应用足够灵活，能够容纳随之产生的任何新代码。

\myGraphic{0.786}{images/CH03_F03_UN07_Mak.png}{}

\begin{mySidebar}{软件工程}

需求收集与分析是软件工程这一更广泛主题中的一个重要课题。软件工程涵盖项目层面的诸多主题，包括项目管理与进度安排；开发方法学（如极限编程和测试驱动开发）；企业应用架构；任务管理、缺陷跟踪和软件版本控制工具；测试策略；项目文档；产品部署与维护；以及其他主题。除软件设计这一主题外，软件工程总体上不在本书的讨论范围之内。

关于如何收集和分析软件需求，已经有整本的专著。本书只能对这一过程做一个简要概述。
\end{mySidebar}

\section{用于创建和记录设计的统一建模语言图}

统一建模语言（Unified Modeling Language，UML）是一套业界标准的图表族，用于帮助软件开发者设计应用并记录设计。下一节将介绍 UML 用例图，它向应用的客户和开发者描述最终用户将如何与应用交互。第 4 章将介绍 UML 类图、状态图和序列图，它们进一步记录应用的设计。

UML 绘图是每个软件开发者技能库中一项有用的技能。这些图不仅对记录应用的设计很重要，其直观的可视化特性也使其在软件设计和编码过程中非常有帮助。

\begin{myWarning}{注意}
大多数流行的计算机绘图工具都提供 UML 对象的面板。在开发迭代过程中，这些图易于创建和操作，便于跟踪应用不断演进的设计。
\end{myWarning}

\section{用例为需求提供上下文}

应用的功能性需求和非功能性需求清单，是开发团队为描述应用而编写的文档的一部分。用例为需求提供上下文，并展示需求如何决定应用的运行时行为。

用例描述应用的最终用户（通常具有某种特定角色，例如图书目录应用中的顾客或图书管理员）如何执行一系列操作来实现某个特定目标。因此，我们从该用户的视角来编写用例。由于用户可以用不同的方式与应用交互以实现各种目标，我们需要多个用例来涵盖主要目标。

用例文档通常由两部分组成。第一部分是 UML \emph{用例图}，它把若干相关的用例归为一组，并展示哪些用户与哪些用例交互。第二部分是为每个用例编写的\emph{用例描述}。

\subsection{3.7.1 UML 用例图}

图 3.4 展示了一个图书目录应用高级版本的示例用例图。在这个版本中，后端数据库存储图书及其属性。

\myGraphic{0.739}{images/CH03_F04_Mak.png}{图 3.4 一个 UML 用例图，包含图书目录应用高级版本中六个用例的名称。虚线框中的标签和虚线箭头标示了该图的各个组成部分。绘制用例图时，不要包含这些标签和箭头。}

图中展示了若干用例的名称，以及不同的参与者如何与它们交互。包围用例的方框代表应用及其边界——了解应用中包含什么、不包含什么很重要。每个用例的名称都应简短、具有描述性，并采用“动词—名词”形式，例如“Add Book”和“Log in customer”。人形图标代表与用例交互的参与者。一个\emph{参与者}（actor）是应用外部任何与应用交互的实体，它可以是一个人（例如图书管理员和顾客），也可以是另一个应用或系统（例如数据库系统）。交互线表示每个参与者可以与哪些用例交互。

\myGraphic{0.651}{images/CH03_F04_UN08_Mak.png}{}

大型应用可能有许多用例图，每张图包含少量（最多大约六个）用例。

\subsection{3.7.2 用例描述}

用例为应用的需求提供上下文。它们展示功能性需求如何决定应用的功能，也展示哪些非功能性需求能确保这些功能切实可用。

我们为用例图中的每个用例单独编写一份用例描述。该描述提供有关该用例的详细信息：

\begin{itemize}
\item \emph{用例名称}——应简短、具有描述性，并采用“动词—名词”形式。
\item \emph{目标}——参与者试图实现什么？
\item \emph{用例摘要}——一段简短（一到两句话）的描述。
\item \emph{参与者}——谁或什么与该用例交互？
\item \emph{前置条件}——在该用例开始执行之前，必须满足什么条件，或必须已经发生了什么？这里可以引用其他用例。
\item \emph{触发器}——参与者做了什么来启动该用例？请从该参与者的视角编写下面的一系列操作步骤。
\item \emph{主操作序列}——在该用例中，触发它的参与者（以及其他任何参与者）与应用之间发生的一系列操作步骤。步骤不应超过大约 10 个。如果需要更多步骤，该用例可能过于复杂，应当加以拆分。这些步骤可以引用其他用例。
\item \emph{备用操作序列}——如果主序列中出了差错，应该发生哪些操作步骤？
\item \emph{后置条件}——该用例结束后将处于什么状态？
\item \emph{非功能性需求}——哪些非功能性需求适用于该用例？
\item \emph{术语表}——定义该用例中读者可能感到困惑的任何术语。
\end{itemize}

从触发该用例的参与者的视角来编写每个用例。

下面是图书目录应用高级版本的一个用例描述示例。顾客填写一个包含“Search”按钮的在线表单，以指定想要查询的图书目标属性，从而进行图书搜索。

\begin{itemize}
\item \emph{用例名称}——Search Catalogue。
\item \emph{目标}——在图书目录中查找与顾客目标属性匹配的图书。
\item \emph{用例摘要}——顾客使用一组图书目标属性搜索图书目录，目录返回所有与这些属性匹配的图书清单。
\item \emph{参与者}——顾客和后端数据库。
\item \emph{前置条件}——图书及其属性已加载到目录中。参见用例 Add Book。顾客已完成填写 Attributes 表单，填入了目标图书属性。参见用例 Complete Form。
\item \emph{触发器}——顾客按下“Search”按钮。
\item \emph{主操作序列}： 1. 应用验证 Attributes 表单已正确填写。 2. 应用根据表单中的图书属性生成一条数据库查询。 3. 应用把查询发送给后端数据库服务器。 4. 数据库向应用返回匹配图书的清单。 5. 应用对匹配图书的清单进行格式化以便展示。 6. 顾客看到匹配图书的清单。
\item \emph{备用操作序列 1}——表单填写不正确。替换主序列的步骤 2–6。 2. 高亮显示填写错误的表单字段。 3. 显示一条说明性的错误消息。 4. 顾客更正出错的表单字段。 5. 返回主序列的步骤 1。
\item \emph{备用操作序列 2}——没有匹配的图书。替换主序列的步骤 4–6。 4. 数据库返回一个空清单。 5. 应用显示“No books found.”。
\item \emph{后置条件}——顾客看到一份匹配图书的清单或者“No books found.”这条消息。图书目录没有任何变化。
\item \emph{非功能性需求}： 搜索结果必须在 2 秒以内返回。 顾客应处于 Windows、macOS 或 Linux 平台上。 应用必须能被母语为英语、西班牙语或越南语的顾客使用。
\item \emph{术语表}： \emph{Catalogue}——可搜索的图书及其属性仓库 \emph{Attribute}——图书的一项特征，可在搜索时进行匹配，例如书名或作者姓名 \emph{Customer}——搜索图书目录的用户
\end{itemize}

用例描述应当简短、简单、不拘形式。应用的客户和开发者都需要能够理解它们。因此，我们不应使用过于技术化的语言。描述中不应包含实现细节。我们必须专注于应用在响应参与者的触发动作时需要做什么，而不是应用将如何去做。

需求来自与应用客户持续不断的沟通协作。把目前已知的需求写下来，并创建用例。用例向客户展示应用在各种场景下将如何表现。这能激发更多需求。新的需求又可能需要新的用例。可能会经历好几轮“新增需求、创建新用例”的过程。一个具有丰富参与者交互的大型应用可能有几十个用例。

引出需求的一个有力方法是把应用的原型展示给客户看。这个原型可以是一些粗糙但能用的代码，用来演示几个关键用例的实际运行。看到一段实实在在能运行的代码，哪怕拼凑得很简陋，往往也能激发客户指出哪些地方缺失、哪些多余，或者哪些没有按预期工作。这个原型也可以简单到只是一套幻灯片，里面包含应用使用过程的模拟截图。然后我们可以把客户的评论和观察结果整理成新的需求，或者据此修改已有的需求。

\section{功能规格说明与软件确认}

开发团队编写\emph{功能规格说明}（functional specification），用来记录即将创建的应用。它的目的是用非技术性、不含术语的语言，向应用的客户和开发者双方说明情况。

\begin{myWarning}{注意}
功能规格说明的格式和内容通常由客户所在的组织决定。不同的组织可能给这份文档起不同的名称，例如 External Reference Specification。之所以说“外部”，是因为这份文档从外部视角看待应用。它不应包含任何内部实现细节。实现细节属于第 4 章介绍的设计规格说明。
\end{myWarning}

功能规格说明应包含以下内容：

\begin{itemize}
\item \emph{应用名称}——例如 Book Catalogue。
\item \emph{清晰的问题陈述}——该应用要解决什么问题？例如，能够存储和搜索图书。
\item \emph{目标}——该应用要达成什么？例如，创建一种把图书录入仓库、然后用目标图书属性搜索它们的途径。
\item \emph{功能性需求}——功能性需求的清单。需求应使用 \emph{must} 或 \emph{shall} 这类情态动词以强烈的语气表述。
\item \emph{非功能性需求}——非功能性需求的清单。需求应使用 \emph{must} 或 \emph{shall} 这类情态动词以强烈的语气表述。
\item \emph{用例}——UML 用例图，以及每个用例的用例描述。
\end{itemize}

功能规格说明还可以包括外部测试计划、部署计划和维护计划。这些通常是独立的文档。

如果包含外部测试计划，它应描述黑盒测试。这类测试无需了解应用代码的内部实现即可进行。对于每项测试，测试计划会描述输入数据或用户操作应该是什么，以及预期的结果。黑盒测试有助于验证应用是否满足其功能性和非功能性需求。这类测试通常由不属于开发团队的测试工程师来执行。

\myGraphic{0.698}{images/CH03_F04_UN09_Mak.png}{}

功能规格说明应该写得多详尽？这往往取决于客户所在的组织和开发方法学。有些组织和方法学不太强调文档，更青睐原型开发以及与客户更紧密的协作。另一些组织则认为功能规格说明是一份极其重要的文档，甚至走极端，把它当作客户与软件开发者之间的合同。它们可能要求多个层级的管理者签字批准一份完整的功能规格说明，然后认定其已冻结，才允许开发者着手设计应用。

但理想情况下，正如 图 3.1 中虚线时间段所暗示的，当我们在应用设计和编码过程中需要微调需求或发现新需求时，应当获准修改功能规格说明。它应当在某个时间点之前都可修改，这个时间点通常称为“功能冻结”（features freeze），由开发者和客户共同商定。尽管最初可以变动，功能规格说明对开发者和客户双方而言仍是一个重要的基准。

无论它是否上升到合同的高度，应用客户仔细阅读功能规格说明都至关重要，而该文档不应包含实现细节。通过阅读、理解并批准功能规格说明，客户\emph{确认}（validate）了应用，也就是确认我们即将构建的是正确的应用，一个满足客户需求的应用。

一旦应用得到确认（记住，客户仍可能改变对需求的想法），我们就可以开始“设计—编码—测试”的迭代。测试是对我们是否正确地构建软件的\emph{验证}（verification）——验证我们正在实现无缺陷的代码，以满足每一条需求。

\begin{myWarning}{注意}
软件验证与确认（software verification and validation，software V\&V）是软件工程和软件质量控制领域中的一个重要课题。
\end{myWarning}

\section{类从何而来？}

一旦我们至少有了功能规格说明的初稿，就可以开始考虑应用的类了。要创建初始的类集合，需要分析应用的功能性需求。需求分析（在设计中仔细审查应用需求的过程）是软件工程的一个重要课题。本章余下的部分将介绍其中最重要的要点。

\myGraphic{0.980}{images/CH03_F04_UN10_Mak.png}{}

下面是第 2 章中那个更高级版本图书目录应用的一组扩展功能性需求。回忆一下，功能性需求规定了应用必须做什么，或必须允许用户做什么：

\begin{itemize}
\item 图书目录应存储不同种类的图书及其属性。
\item 图书管理员必须能够向目录中添加新书。
\item 图书管理员必须能够更新和删除目录中已有的图书。
\item 图书的种类应包括小说、烹饪书和实用指南。
\item 所有图书都必须有书名、作者姓氏和作者名字的属性值。
\item 小说必须包含出版年份和体裁属性。
\item 体裁必须包括冒险、经典、侦探、奇幻、历史、恐怖、言情和科幻。
\item 烹饪书必须包含地区属性。
\item 实用指南类图书必须包含主题属性。
\item 顾客必须能够通过提供任意数量的所需目标属性值来搜索目录。
\item 顾客必须填写一个基于网络浏览器的表单，以指定图书搜索的目标属性值。
\item 顾客在表单中的输入必须经过验证，以确认格式和取值正确。
\item 搜索时，字符串属性的匹配必须不区分大小写。
\item 顾客必须能够指定任意数量的不关心（通配）目标属性。
\item 默认情况下，每个不关心的属性都必须与目录中所有图书的对应属性匹配。
\item 如果目录中的一本书拥有顾客目标属性中的所有属性，且所有对应的图书属性值与目标属性值都相等，或者该目标属性是不关心的属性，则这本书匹配。
\end{itemize}

\subsection{3.9.1 文本分析：名词可以变成类}

回忆 1.8 节，类规定了其对象在运行时的状态和行为。要确定我们的应用应该有哪些类，首先找出需求中的名词：\emph{图书}（book）、\emph{目录}（catalogue）、\emph{属性}（attribute）、\emph{图书管理员}（librarian）、\emph{顾客}（customer）、\emph{种类}（kind）、\emph{书名}（title）、\emph{姓名}（name）、\emph{年份}（year）、\emph{体裁}（genre）、\emph{地区}（region）、\emph{主题}（subject）、\emph{浏览器}（browser）、\emph{表单}（form）、\emph{输入}（input）和 \emph{字符串}（string）。这些名词代表了潜在的类，但如表 3.1 所示，哪些名词应该成为类、哪些应该成为应用中的属性，需要靠判断来决定。这类设计决策可能需要在开发迭代中反复斟酌。

\begin{center}
{\bfseries 表 3.1 判断需求中的每个名词是否应成为一个类。}\par\vspace{3bp}
\begin{tabularx}{\textwidth}{XX}
\toprule
名词 & 是类吗？ \\
\midrule
catalogue & 是。应用实现了一个图书目录。 \\
Book & 是。目录存储图书对象。 \\
attribute & 是。每个图书对象都有一组属性。 \\
librarian & 否。根据图 3.4，图书管理员是应用外部的实体。 \\
customer & 否。根据图 3.4，顾客是应用外部的实体。 \\
kind & 否。它是一个枚举常量形式的属性值。 \\
title & 否。它是一个字符串形式的属性值。 \\
name & 否。它是一个字符串形式的属性值。 \\
year & 否。它是一个整数形式的属性值。 \\
genre & 否。它是一个枚举常量形式的属性值。 \\
region & 否。它是一个枚举常量形式的属性值。 \\
subject & 否。它是一个枚举常量形式的属性值。 \\
browser & 否。应用与现有的浏览器协同工作。 \\
form & 是。应用管理一个用户输入表单。 \\
input & 否。它是用户输入到表单中的属性值。 \\
format & 否。格式在应用中并不是一个独立的事物。 \\
value & 否。属性值是整数、字符串或枚举常量。 \\
string & 否。应用使用 C++ 内置的字符串。 \\
\bottomrule
\end{tabularx}
\end{center}

剩下这几个可以成为应用初始类的名词：\emph{catalogue}、\emph{book}、\emph{attribute} 和 \emph{form}。对于每个类，我们必须确定它的成员变量，以便其对象能在运行时维护状态。表 3.2 可以是像第 2 章所述那样经过若干次设计迭代后的结果。后续的迭代还可能发现更多的类。

\begin{center}
{\bfseries 表 3.2 应用的初始类}\par\vspace{3bp}
\begin{tabularx}{\textwidth}{XXX}
\toprule
类 & 状态 & 成员变量 \\
\midrule
\texttt{Catalogue} & 图书清单 & \texttt{booklist}（指向 \texttt{Book} 对象的指针的 vector） \\
\texttt{Book} & 图书属性 & \texttt{attributes}（指向该书属性的指针） \\
\texttt{Attributes} & 属性值 & \texttt{attribute\_\allowbreak{}}\texttt{map}（键—值对的 map） \\
\texttt{Form} & 输入值 & 为各个图书属性分别设置的成员变量 \\
\bottomrule
\end{tabularx}
\end{center}

\subsection{3.9.2 文本分析：动词可以变成成员函数}

接下来，我们必须根据需求中的动词来确定每个类的对象在运行时的行为：\emph{store}、\emph{add}、\emph{update}、\emph{delete}、\emph{include}、\emph{have}、\emph{be}、\emph{search}、\emph{complete}、\emph{verify}、\emph{constitute}、\emph{specify} 和 \emph{match}。我们只应考虑那些对对象执行某种动作的（及物）动词。与表 3.2 一样，表 3.3 可以是经过若干次设计迭代后的结果。

\begin{center}
{\bfseries 表 3.3 确定每个类的行为。}\par\vspace{3bp}
\begin{tabularx}{\textwidth}{XXX}
\toprule
动词 & 类 & 成员函数 \\
\midrule
add & \texttt{Catalogue} & \texttt{add()} 向 booklist 添加一本书 \\
update & \texttt{Catalogue} & \texttt{update()} 更新 booklist 中的一本书 \\
delete & \texttt{Catalogue} & \texttt{delete()} 从 booklist 中删除一本书 \\
search & \texttt{Catalogue} & \texttt{find()} 在 booklist 中查找匹配的图书 \\
verify & \texttt{Form} & \texttt{verify()} 验证表单字段中的用户输入 \\
match & \texttt{Attributes} & \texttt{is\_\allowbreak{}match(\allowbreak{})\allowbreak{}} 检查属性对是否匹配 \\
\bottomrule
\end{tabularx}
\end{center}

每个类的成员函数都操作该类的成员变量。例如，在类 \texttt{Catalogue} 中，成员函数 \texttt{add()} 把指向 \texttt{Book} 对象的指针添加到成员变量 \texttt{booklist} 中。

\myGraphic{0.809}{images/CH03_F04_UN11_Mak.png}{}

现在我们已经有了初步的类，接下来要确保把它们以及后续出现的所有类都设计好。这就是接下来几章的主题。

\section*{小结}

\begin{itemize}
\item 功能性需求规定了应用为取得成功必须做什么，或必须允许用户做什么。
\item 非功能性需求规定了应用为取得成功必须满足的限制或约束。它们与功能性需求同等重要。例如性能需求或平台需求。
\item 使用 \emph{must} 和 \emph{shall} 这类语气强烈的动词来表述需求。
\item 获取好的需求至关重要。访谈应用的客户，或者观察人们当前如何完成任务。制作应用的原型并展示给客户看。设想应用可以怎样提高他们的生产力。
\item 需求来自与客户的持续协作，而客户并不总是知道自己想要什么。要做好准备，需求可能在应用设计期间、甚至编码期间发生变化。
\item 好的需求是清晰、一致、正确、完整、切合实际、可验证且可追溯的。
\item 用例为需求提供上下文，并展示需求如何决定应用的功能。用 UML 用例图和用例描述来记录用例。用例可以激发新的需求。
\item 用非技术性、不含术语的语言编写的功能规格说明，向客户和软件开发者双方说明应用将做什么。它可以包含用于黑盒测试的外部测试计划。有些组织把这份文档视为客户与开发者之间的合同。理想情况下，它应当可修改（直到功能冻结这一里程碑为止），以便适应需求的变化。
\item 客户通过批准功能规格说明来确认（validate）软件开发者将构建正确的应用。开发者通过测试来验证（verify）他们正确地构建了应用。软件验证与确认（software V\&V）是软件工程和软件质量控制领域中的一个重要课题。
\item 分析需求，以得到应用的初始类集合。确定哪些名词成为类、哪些动词成为实现行为的成员函数。为各个类分派成员变量和成员函数。
\item 不要陷入“分析瘫痪”。制作并展示原型，开始开发迭代，更多的类自然会浮现出来。
\end{itemize}
